\DocumentMetadata{
  pdfversion=2.0,
  pdfstandard=ua-2,
  lang=en-US,
  tagging=on,
  tagging-setup={math/setup=mathml-SE,tabsorder=structure}
}
\documentclass[11pt,letterpaper]{article}
\usepackage[margin=0.68in,includefoot]{geometry}
\usepackage{newcomputermodern}
\usepackage{amsmath,amscd,mathtools}
\usepackage{tikz,adjustbox}
\providecommand{\square}{\mdlgwhtsquare}
\usepackage[hidelinks]{hyperref}
\hypersetup{
  pdftitle={Numerical Analysis PhD Exam, January 2025},
  pdfauthor={University of Florida Department of Mathematics},
  pdfsubject={Graduate mathematics examination},
  pdfdisplaydoctitle=true,
  bookmarksopen=true
}
\setcounter{secnumdepth}{0}
\setlength{\parindent}{0pt}
\setlength{\parskip}{0.28em}
\setlength{\emergencystretch}{3em}
\setlength{\leftmargini}{2.15em}
\setlength{\leftmarginii}{2.65em}
\setlength{\leftmarginiii}{3em}
\pagestyle{empty}
\raggedbottom
\allowdisplaybreaks
\NewTaggingSocketPlug{block/list/label}{examproblem}{%
  \tagstructbegin{tag=Lbl}\tagstructend%
  \tagstructbegin{tag=\LBody}%
  \tagstructbegin{tag=\ProblemHeadingTag,title-o={\ProblemHeadingText}}%
  \tagmcbegin{tag=\ProblemHeadingTag,actualtext={\ProblemHeadingText}}%
  #2%
  \tagmcend%
  \tagstructend%
}
\newenvironment{examproblems}[2]{%
  \def\ProblemHeadingTag{#2}%
  \AssignTaggingSocketPlug{block/list/label}{examproblem}%
  \begin{enumerate}%
  \setcounter{enumi}{#1}%
  \renewcommand{\labelenumi}{\textbf{\theenumi.}}%
  \setlength{\topsep}{0.35em}%
  \setlength{\partopsep}{0pt}%
  \setlength{\itemsep}{0.48em}%
  \setlength{\parsep}{0.18em}%
}{\end{enumerate}}
\newenvironment{examsubparts}{%
  \AssignTaggingSocketPlug{block/list/label}{default}%
  \begin{enumerate}%
  \setlength{\topsep}{0.18em}%
  \setlength{\partopsep}{0pt}%
  \setlength{\itemsep}{0.16em}%
  \setlength{\parsep}{0.1em}%
}{\end{enumerate}}
\newenvironment{examinstructions}{%
  \par\begingroup\itshape
}{%
  \par\endgroup\vspace{0.18em}
}
\makeatletter
\renewcommand\subsection{\@startsection{subsection}{2}{0pt}%
  {-1.25ex plus -.25ex minus -.1ex}%
  {0.55ex plus .1ex}%
  {\normalfont\large\bfseries}}
\renewcommand{\@maketitle}{%
  \newpage\null\vskip -1em%
  {\footnotesize\bfseries
    \href{https://www.ufl.edu/}{University of Florida}, \href{https://math.ufl.edu/}{Department of Mathematics}\par}%
  \vskip -0.5em%
  \tagstructbegin{tag=H1,title={Numerical Analysis PhD Exam, January 2025}}%
  \tagpdfparaOff%
  \tagmcbegin{tag=H1}%
    {\LARGE\bfseries\href{https://gma.math.ufl.edu/exams/numerical-analysis-phd/}{Numerical Analysis PhD Exam}, January 2025\par}%
  \tagmcend%
  \tagpdfparaOn%
  \tagstructend%
  \vskip 0.25em%
  \tagmcbegin{artifact}%
    \hrule height 1pt%
  \tagmcend%
  \vskip 1em%
}
\makeatother
\title{Numerical Analysis PhD Exam, January 2025}
\author{}
\date{}
\begin{document}
\maketitle
\thispagestyle{empty}
\begin{examinstructions}
Do 4 (four) of the first 5 (1-5) and 4 (four) of the last 5 problems (6-10).
\end{examinstructions}
\begin{examproblems}{0}{H2}
\def\ProblemHeadingText{Problem 1.}
\item
Assume \(A\in\mathbb{C}^{m\times m}\).
\begin{examsubparts}
\item[\textnormal{(a)}]
Show that~\(A\) has a Schur decomposition.
\item[\textnormal{(b)}]
If~\(A\) has a collection of~\(m\) linearly independent eigenvectors, show that~\(A\) is diagonalizable.
\end{examsubparts}
\def\ProblemHeadingText{Problem 2.}
\item
If \(q_1,\ldots,q_n\) is an orthonormal basis for the subspace \(V\subset\mathbb{C}^m\) with \(m>n\), prove that the orthogonal projector onto~\(V\) is \(QQ^*\), where~\(Q\) is the matrix whose columns are the \(q_j\).
\def\ProblemHeadingText{Problem 3.}
\item
Let \(A\in\mathbb{C}^{m\times n}\), with \(m\geq n\) and \(\operatorname{rank}(A)=n\geq 3\). Let \(a_1,a_2,\ldots\) denote the columns of~\(A\).
\begin{examsubparts}
\item[\textnormal{(a)}]
Using the modified Gram–Schmidt process, write out expressions for \(q_1,q_2,q_3\), the first three columns of~\(Q\) in the QR decomposition of~\(A\).
\item[\textnormal{(b)}]
Show the vector \(q_3\) found in part (a) is orthogonal to both \(q_1\) and \(q_2\).
\end{examsubparts}
\def\ProblemHeadingText{Problem 4.}
\item
Let \(\lVert\cdot\rVert\) be a subordinate (induced) matrix norm.
\begin{examsubparts}
\item[\textnormal{(a)}]
If~\(E\) is a \(n\times n\) with \(\lVert E\rVert<1\), then show \(I+E\) is nonsingular and
\[
\lVert(I+E)^{-1}\rVert\leq\frac{1}{1-\lVert E\rVert}.
\]
\item[\textnormal{(b)}]
If~\(A\) is a \(n\times n\) invertible and~\(E\) is \(n\times n\) with \(\lVert A^{-1}\rVert\lVert E\rVert<1\), then show \(A+E\) is nonsingular and
\[
\lVert(A+E)^{-1}\rVert\leq\frac{\lVert A^{-1}\rVert}{1-\lVert A^{-1}\rVert\lVert E\rVert}.
\]
\end{examsubparts}
\def\ProblemHeadingText{Problem 5.}
\item
Let \(A=U\Sigma V^*\) be the singular value decomposition of \(A\in\mathbb{C}^{m\times n}\). Let \(u_j\) denote column~\(j\) of~\(U\).
\begin{examsubparts}
\item[\textnormal{(a)}]
Suppose \(\operatorname{rank}(A)=p<n<m\). Show \(\{u_1,u_2,\ldots,u_p\}\) is a basis for \(\operatorname{Col}(A)\) and \(\{u_{p+1},u_{p+2},\ldots,u_m\}\) is a basis for \(\operatorname{Null}(A^*)\).
\item[\textnormal{(b)}]
Suppose~\(A\) is full rank and \(x\neq 0\). Let \(\sigma_i\), \(i=1,\ldots,n\), be the singular values of~\(A\). Show
\[
\sigma_1\geq\frac{\lVert Ax\rVert_2}{\lVert x\rVert_2}\geq\sigma_n>0.
\]
 If you want to use the property that \(\lVert A\rVert_2=\sigma_1\), then you must prove that it holds.
\end{examsubparts}
\def\ProblemHeadingText{Problem 6.}
\item
\(g(x)=\frac{x}{2}+\frac{a}{2x}\) has \(x=\sqrt{a}\) as a fixed point where \(a>0\). Suppose we use the fixed point method with a starting point \(x_0>\sqrt{a}\). Use the theory of the fixed point method to determine the condition on \(x_0\) which guarantees
\[
|x_{k+1}-\sqrt{a}|\leq\frac{1}{2}|x_k-\sqrt{a}|,\qquad k\geq 1.
\]
\def\ProblemHeadingText{Problem 7.}
\item
Let \(f\in C[0,1]\) and let \(\sum_{i=1}^n w_i f(x_i)\) (\(n\geq 1\)) be an approximation to \(\int_0^1 f(x)\,dx\). Assume that \(0\leq w_i\leq 1\), \(0\leq x_i\leq 1\) (\(i=1,2,\ldots,n\)). Let
\[
E_n(f)=\int_0^1 f(x)\,dx-\sum_{i=1}^n w_i f(x_i)
\]
 be the error in the approximation. Assume that \(E_n(p_n)=0\) for \(p_n\), which denotes any polynomial of degree \(\leq n\). Use the Weierstrass approximation theorem to prove that, given \(\epsilon>0\), there is an \(N>0\) such that \(|E_n(f)|<\epsilon\) when \(n>N\).
\def\ProblemHeadingText{Problem 8.}
\item
Given a differentiable function \(f(x)\), consider the problem of finding a polynomial \(p(x)\in\mathbb{P}^n\) such that
\[
p(x_0)=f(x_0),\qquad p'(x_i)=f'(x_i),\qquad i=1,2,\ldots,n,
\]
 where \(x_i\), \(i=1,2,\ldots,n\), are distinct nodes. (It is not excluded that \(x_1=x_0\).) Show that the problem has a unique solution and describe how it can be obtained.
\def\ProblemHeadingText{Problem 9.}
\item
Let \(\{p_i(x)\}_{i=0}^{\infty}\) be the sequence of orthogonal polynomials obtained from the Gram–Schmidt orthogonalization process of \(\{x^i\}_{i=0}^{\infty}\). Consider the quadrature formula \(Q(f)=\sum_{j=0}^m\alpha_j f(x_j)\) for \(J(f)=\int_a^b f(x)\rho(x)\,dx\). Prove that \(Q(f)\) is exact for polynomials of degree \(\leq 2m+1\) if and only if the quadrature nodes \(x_j\) (\(j=0,1,\ldots,m\)) are the roots of \(p_{m+1}(x)\) and the quadrature weights \(\alpha_j=\int_a^b l_j(x)\rho(x)\,dx\), (\(j=0,1,\ldots,m\)).
\def\ProblemHeadingText{Problem 10.}
\item
Consider numerically solving the initial value problem
\[
y'(t)=f(t,y),\qquad 0<t\leq t_f,\qquad\text{with }y(0)=\eta.
\]
 Assume~\(f\) is sufficiently differentiable and let~\(h\) denote the step size. Show that the method
\[
y_{n+2}-y_{n+1}=\frac{1}{4}h(f_{n+2}+2f_{n+1}+f_n)
\]
 is \(A_0\) stable.
\end{examproblems}
\end{document}
